\documentclass[10pt]{article}
\usepackage[margin=0.7in]{geometry}
\usepackage{booktabs}
\usepackage{longtable}
\usepackage{lmodern}
\title{HumanEval TRAIN-set typicality metrics --- per-epoch dynamics\\ \large gemma-4-31b $\times$ s2 (RankAlign) \& s4 (self-tc) $\times$ base/ep0/ep1/ep2}
\author{}
\date{2026-06-19}
\begin{document}
\maketitle

One table per metric (gen ROC, val ROC, val acc, Pearson). Each shows the \textbf{raw} and
\textbf{tc} (typicality-corrected) gen variants, with an \textbf{epoch axis}
(base $\to$ ep0 $\to$ ep1 $\to$ ep2) so the training dynamics are visible. Systems:
gemma-4-31b, settings s2 (RankAlign) and s4 (self-tc / New+fsx+self-tc), on both
\texttt{humaneval-v2.1correct-upper} and \texttt{-multi}. Values $\times 100$, 1 decimal;
means over the 82 problems. \textbf{Scores are on the \texttt{--train} split} (the candidate
pool the models were trained against), subsampled to N=50 \emph{stratified} candidates per
problem (25 positive / 25 negative). Eval = (typicality mode)/(base-typ vs no-base).

\medskip\noindent\textit{Notes:} (i) \texttt{val ROC} and \texttt{val acc} are validator-based,
so their raw and tc columns are identical by construction (shown for uniformity). (ii) \texttt{raw}
gen scores do not depend on the eval mode, so the raw column repeats within a (dataset, setting,
epoch) block. (iii) The \texttt{base} row is the untrained gemma-4-31b-it model (shared start point
for both settings). (iv) s4 (self-tc) was evaluated in \emph{self}-typicality mode only
(self/base-typ, self/no-base); base and s2 were evaluated in all four modes. (v) These are
TRAIN-split numbers --- generator/validator agreement is expected to be high on data the model was
trained on; the interest is the \emph{dynamics} (how raw vs tc move across epochs), not the absolute
level. The companion TEST-split tables are in \texttt{docs/he\_harvest\_2026-06-17/}.

{\small
\begin{longtable}{llll rr}
\caption{HumanEval TRAIN-set gen ROC ($\times100$) per-epoch dynamics, raw vs tc gen variant. gemma-4-31b, s2 (RankAlign) \& s4 (self-tc), base/ep0/ep1/ep2, both datasets. Means over 82 problems, N=50 stratified candidates each.}\label{tab:he_train_gen_roc}\\
\toprule
Data & Setting & Epoch & Eval & raw & tc \\
\midrule
\endfirsthead
\multicolumn{6}{l}{\itshape (continued)}\\
\toprule
Data & Setting & Epoch & Eval & raw & tc \\
\midrule
\endhead
\bottomrule
\endfoot
upper & base & base & self/base-typ & 74.1 & 68.8 \\
upper & base & base & self/no-base & 74.1 & 68.8 \\
upper & base & base & neg/base-typ & 74.1 & 72.2 \\
upper & base & base & neg/no-base & 74.1 & 72.2 \\
upper & s2 & ep0 & self/base-typ & 96.5 & 81.6 \\
upper & s2 & ep0 & self/no-base & 96.5 & 85.3 \\
upper & s2 & ep0 & neg/base-typ & 96.5 & 85.9 \\
upper & s2 & ep0 & neg/no-base & 96.5 & 87.5 \\
upper & s2 & ep1 & self/base-typ & 94.4 & 80.0 \\
upper & s2 & ep1 & self/no-base & 94.4 & 82.7 \\
upper & s2 & ep1 & neg/base-typ & 94.4 & 85.4 \\
upper & s2 & ep1 & neg/no-base & 94.4 & 80.8 \\
upper & s2 & ep2 & self/base-typ & 96.3 & 79.4 \\
upper & s2 & ep2 & self/no-base & 96.3 & 83.4 \\
upper & s2 & ep2 & neg/base-typ & 96.3 & 83.4 \\
upper & s2 & ep2 & neg/no-base & 96.3 & 87.2 \\
upper & s4 & ep0 & self/base-typ & 93.1 & 87.4 \\
upper & s4 & ep0 & self/no-base & 93.1 & 83.5 \\
upper & s4 & ep1 & self/base-typ & 91.7 & 90.9 \\
upper & s4 & ep1 & self/no-base & 91.7 & 84.5 \\
upper & s4 & ep2 & self/base-typ & 91.0 & 90.9 \\
upper & s4 & ep2 & self/no-base & 91.0 & 86.1 \\
\midrule
multi & base & base & self/base-typ & 53.0 & 64.6 \\
multi & base & base & self/no-base & 53.0 & 64.6 \\
multi & base & base & neg/base-typ & 53.0 & 66.9 \\
multi & base & base & neg/no-base & 53.0 & 66.9 \\
multi & s2 & ep0 & self/base-typ & 98.9 & 95.4 \\
multi & s2 & ep0 & self/no-base & 98.9 & 95.8 \\
multi & s2 & ep0 & neg/base-typ & 98.9 & 95.7 \\
multi & s2 & ep0 & neg/no-base & 98.9 & 84.3 \\
multi & s2 & ep1 & self/base-typ & 98.1 & 92.0 \\
multi & s2 & ep1 & self/no-base & 98.1 & 94.1 \\
multi & s2 & ep1 & neg/base-typ & 98.1 & 94.7 \\
multi & s2 & ep1 & neg/no-base & 98.1 & 88.5 \\
multi & s2 & ep2 & self/base-typ & 94.6 & 95.2 \\
multi & s2 & ep2 & self/no-base & 94.6 & 90.4 \\
multi & s2 & ep2 & neg/base-typ & 94.6 & 94.6 \\
multi & s2 & ep2 & neg/no-base & 94.6 & 78.7 \\
multi & s4 & ep0 & self/base-typ & 85.1 & 97.9 \\
multi & s4 & ep0 & self/no-base & 85.1 & 89.1 \\
multi & s4 & ep1 & self/base-typ & 86.4 & 98.2 \\
multi & s4 & ep1 & self/no-base & 86.4 & 89.9 \\
multi & s4 & ep2 & self/base-typ & 85.6 & 99.0 \\
multi & s4 & ep2 & self/no-base & 85.6 & 94.2 \\
\bottomrule
\end{longtable}
}

{\small
\begin{longtable}{llll rr}
\caption{HumanEval TRAIN-set val ROC ($\times100$) per-epoch dynamics, raw vs tc gen variant. gemma-4-31b, s2 (RankAlign) \& s4 (self-tc), base/ep0/ep1/ep2, both datasets. Means over 82 problems, N=50 stratified candidates each.}\label{tab:he_train_val_roc}\\
\toprule
Data & Setting & Epoch & Eval & raw & tc \\
\midrule
\endfirsthead
\multicolumn{6}{l}{\itshape (continued)}\\
\toprule
Data & Setting & Epoch & Eval & raw & tc \\
\midrule
\endhead
\bottomrule
\endfoot
upper & base & base & self/base-typ & 90.6 & 90.6 \\
upper & base & base & self/no-base & 90.6 & 90.6 \\
upper & base & base & neg/base-typ & 90.6 & 90.6 \\
upper & base & base & neg/no-base & 90.6 & 90.6 \\
upper & s2 & ep0 & self/base-typ & 92.6 & 92.6 \\
upper & s2 & ep0 & self/no-base & 92.6 & 92.6 \\
upper & s2 & ep0 & neg/base-typ & 92.6 & 92.6 \\
upper & s2 & ep0 & neg/no-base & 92.6 & 92.6 \\
upper & s2 & ep1 & self/base-typ & 90.9 & 90.9 \\
upper & s2 & ep1 & self/no-base & 90.9 & 90.9 \\
upper & s2 & ep1 & neg/base-typ & 90.9 & 90.9 \\
upper & s2 & ep1 & neg/no-base & 90.9 & 90.9 \\
upper & s2 & ep2 & self/base-typ & 90.2 & 90.2 \\
upper & s2 & ep2 & self/no-base & 90.2 & 90.2 \\
upper & s2 & ep2 & neg/base-typ & 90.2 & 90.2 \\
upper & s2 & ep2 & neg/no-base & 90.2 & 90.2 \\
upper & s4 & ep0 & self/base-typ & 95.5 & 95.5 \\
upper & s4 & ep0 & self/no-base & 95.5 & 95.5 \\
upper & s4 & ep1 & self/base-typ & 95.5 & 95.5 \\
upper & s4 & ep1 & self/no-base & 95.5 & 95.5 \\
upper & s4 & ep2 & self/base-typ & 96.2 & 96.2 \\
upper & s4 & ep2 & self/no-base & 96.2 & 96.2 \\
\midrule
multi & base & base & self/base-typ & 88.6 & 88.6 \\
multi & base & base & self/no-base & 88.6 & 88.6 \\
multi & base & base & neg/base-typ & 88.6 & 88.6 \\
multi & base & base & neg/no-base & 88.6 & 88.6 \\
multi & s2 & ep0 & self/base-typ & 90.4 & 90.4 \\
multi & s2 & ep0 & self/no-base & 90.4 & 90.4 \\
multi & s2 & ep0 & neg/base-typ & 90.4 & 90.4 \\
multi & s2 & ep0 & neg/no-base & 90.4 & 90.4 \\
multi & s2 & ep1 & self/base-typ & 92.0 & 92.0 \\
multi & s2 & ep1 & self/no-base & 92.0 & 92.0 \\
multi & s2 & ep1 & neg/base-typ & 92.0 & 92.0 \\
multi & s2 & ep1 & neg/no-base & 92.0 & 92.0 \\
multi & s2 & ep2 & self/base-typ & 94.1 & 94.1 \\
multi & s2 & ep2 & self/no-base & 94.1 & 94.1 \\
multi & s2 & ep2 & neg/base-typ & 94.1 & 94.1 \\
multi & s2 & ep2 & neg/no-base & 94.1 & 94.1 \\
multi & s4 & ep0 & self/base-typ & 98.9 & 98.9 \\
multi & s4 & ep0 & self/no-base & 98.9 & 98.9 \\
multi & s4 & ep1 & self/base-typ & 99.5 & 99.5 \\
multi & s4 & ep1 & self/no-base & 99.5 & 99.5 \\
multi & s4 & ep2 & self/base-typ & 99.8 & 99.8 \\
multi & s4 & ep2 & self/no-base & 99.8 & 99.8 \\
\bottomrule
\end{longtable}
}

{\small
\begin{longtable}{llll rr}
\caption{HumanEval TRAIN-set val acc ($\times100$) per-epoch dynamics, raw vs tc gen variant. gemma-4-31b, s2 (RankAlign) \& s4 (self-tc), base/ep0/ep1/ep2, both datasets. Means over 82 problems, N=50 stratified candidates each.}\label{tab:he_train_val_acc}\\
\toprule
Data & Setting & Epoch & Eval & raw & tc \\
\midrule
\endfirsthead
\multicolumn{6}{l}{\itshape (continued)}\\
\toprule
Data & Setting & Epoch & Eval & raw & tc \\
\midrule
\endhead
\bottomrule
\endfoot
upper & base & base & self/base-typ & 70.0 & 70.0 \\
upper & base & base & self/no-base & 70.0 & 70.0 \\
upper & base & base & neg/base-typ & 70.0 & 70.0 \\
upper & base & base & neg/no-base & 70.0 & 70.0 \\
upper & s2 & ep0 & self/base-typ & 74.0 & 74.0 \\
upper & s2 & ep0 & self/no-base & 74.0 & 74.0 \\
upper & s2 & ep0 & neg/base-typ & 74.0 & 74.0 \\
upper & s2 & ep0 & neg/no-base & 74.0 & 74.0 \\
upper & s2 & ep1 & self/base-typ & 74.0 & 74.0 \\
upper & s2 & ep1 & self/no-base & 74.0 & 74.0 \\
upper & s2 & ep1 & neg/base-typ & 74.0 & 74.0 \\
upper & s2 & ep1 & neg/no-base & 74.0 & 74.0 \\
upper & s2 & ep2 & self/base-typ & 72.0 & 72.0 \\
upper & s2 & ep2 & self/no-base & 72.0 & 72.0 \\
upper & s2 & ep2 & neg/base-typ & 72.0 & 72.0 \\
upper & s2 & ep2 & neg/no-base & 72.0 & 72.0 \\
upper & s4 & ep0 & self/base-typ & 76.0 & 76.0 \\
upper & s4 & ep0 & self/no-base & 76.0 & 76.0 \\
upper & s4 & ep1 & self/base-typ & 80.0 & 80.0 \\
upper & s4 & ep1 & self/no-base & 80.0 & 80.0 \\
upper & s4 & ep2 & self/base-typ & 80.0 & 80.0 \\
upper & s4 & ep2 & self/no-base & 80.0 & 80.0 \\
\midrule
multi & base & base & self/base-typ & 74.0 & 74.0 \\
multi & base & base & self/no-base & 74.0 & 74.0 \\
multi & base & base & neg/base-typ & 74.0 & 74.0 \\
multi & base & base & neg/no-base & 74.0 & 74.0 \\
multi & s2 & ep0 & self/base-typ & 76.0 & 76.0 \\
multi & s2 & ep0 & self/no-base & 76.0 & 76.0 \\
multi & s2 & ep0 & neg/base-typ & 76.0 & 76.0 \\
multi & s2 & ep0 & neg/no-base & 76.0 & 76.0 \\
multi & s2 & ep1 & self/base-typ & 74.0 & 74.0 \\
multi & s2 & ep1 & self/no-base & 74.0 & 74.0 \\
multi & s2 & ep1 & neg/base-typ & 74.0 & 74.0 \\
multi & s2 & ep1 & neg/no-base & 74.0 & 74.0 \\
multi & s2 & ep2 & self/base-typ & 78.0 & 78.0 \\
multi & s2 & ep2 & self/no-base & 78.0 & 78.0 \\
multi & s2 & ep2 & neg/base-typ & 78.0 & 78.0 \\
multi & s2 & ep2 & neg/no-base & 78.0 & 78.0 \\
multi & s4 & ep0 & self/base-typ & 88.0 & 88.0 \\
multi & s4 & ep0 & self/no-base & 88.0 & 88.0 \\
multi & s4 & ep1 & self/base-typ & 80.0 & 80.0 \\
multi & s4 & ep1 & self/no-base & 80.0 & 80.0 \\
multi & s4 & ep2 & self/base-typ & 84.0 & 84.0 \\
multi & s4 & ep2 & self/no-base & 84.0 & 84.0 \\
\bottomrule
\end{longtable}
}

{\small
\begin{longtable}{llll rr}
\caption{HumanEval TRAIN-set Pearson ($\times100$) per-epoch dynamics, raw vs tc gen variant. gemma-4-31b, s2 (RankAlign) \& s4 (self-tc), base/ep0/ep1/ep2, both datasets. Means over 82 problems, N=50 stratified candidates each.}\label{tab:he_train_pearson}\\
\toprule
Data & Setting & Epoch & Eval & raw & tc \\
\midrule
\endfirsthead
\multicolumn{6}{l}{\itshape (continued)}\\
\toprule
Data & Setting & Epoch & Eval & raw & tc \\
\midrule
\endhead
\bottomrule
\endfoot
upper & base & base & self/base-typ & 39.5 & 11.2 \\
upper & base & base & self/no-base & 39.5 & 11.2 \\
upper & base & base & neg/base-typ & 39.5 & 30.0 \\
upper & base & base & neg/no-base & 39.5 & 30.0 \\
upper & s2 & ep0 & self/base-typ & 51.1 & 22.7 \\
upper & s2 & ep0 & self/no-base & 51.1 & 27.2 \\
upper & s2 & ep0 & neg/base-typ & 51.1 & 30.7 \\
upper & s2 & ep0 & neg/no-base & 51.1 & 49.9 \\
upper & s2 & ep1 & self/base-typ & 50.0 & 20.8 \\
upper & s2 & ep1 & self/no-base & 50.0 & 26.2 \\
upper & s2 & ep1 & neg/base-typ & 50.0 & 32.3 \\
upper & s2 & ep1 & neg/no-base & 50.0 & 47.4 \\
upper & s2 & ep2 & self/base-typ & 51.1 & 18.6 \\
upper & s2 & ep2 & self/no-base & 51.1 & 27.4 \\
upper & s2 & ep2 & neg/base-typ & 51.1 & 25.5 \\
upper & s2 & ep2 & neg/no-base & 51.1 & 46.9 \\
upper & s4 & ep0 & self/base-typ & 53.3 & 40.3 \\
upper & s4 & ep0 & self/no-base & 53.3 & 36.3 \\
upper & s4 & ep1 & self/base-typ & 57.2 & 54.4 \\
upper & s4 & ep1 & self/no-base & 57.2 & 45.9 \\
upper & s4 & ep2 & self/base-typ & 57.3 & 57.8 \\
upper & s4 & ep2 & self/no-base & 57.3 & 52.5 \\
\midrule
multi & base & base & self/base-typ & 15.8 & 19.4 \\
multi & base & base & self/no-base & 15.8 & 19.4 \\
multi & base & base & neg/base-typ & 15.8 & 19.4 \\
multi & base & base & neg/no-base & 15.8 & 19.4 \\
multi & s2 & ep0 & self/base-typ & 53.3 & 57.8 \\
multi & s2 & ep0 & self/no-base & 53.3 & 57.1 \\
multi & s2 & ep0 & neg/base-typ & 53.3 & 59.4 \\
multi & s2 & ep0 & neg/no-base & 53.3 & 47.0 \\
multi & s2 & ep1 & self/base-typ & 59.8 & 59.8 \\
multi & s2 & ep1 & self/no-base & 59.8 & 60.8 \\
multi & s2 & ep1 & neg/base-typ & 59.8 & 62.5 \\
multi & s2 & ep1 & neg/no-base & 59.8 & 52.9 \\
multi & s2 & ep2 & self/base-typ & 57.3 & 61.5 \\
multi & s2 & ep2 & self/no-base & 57.3 & 56.3 \\
multi & s2 & ep2 & neg/base-typ & 57.3 & 62.1 \\
multi & s2 & ep2 & neg/no-base & 57.3 & 43.9 \\
multi & s4 & ep0 & self/base-typ & 44.9 & 70.5 \\
multi & s4 & ep0 & self/no-base & 44.9 & 61.2 \\
multi & s4 & ep1 & self/base-typ & 47.4 & 72.1 \\
multi & s4 & ep1 & self/no-base & 47.4 & 60.4 \\
multi & s4 & ep2 & self/base-typ & 53.6 & 79.8 \\
multi & s4 & ep2 & self/no-base & 53.6 & 65.0 \\
\bottomrule
\end{longtable}
}


\end{document}
